% !TEX root = ../double_trouble_main.tex

Before choosing a particular model, we describe a general Bayesian framework for predicting 
follow-up variant numbers given pilot data in two populations.

\revision{We start with intuition for our likelihood model. We assume that samples within a single population are exchangeable; that is,
we assume a distribution over the samples and further that the probability of the data is unchanged 
when the samples are observed in a different order.
Second, we observe that every existing work in a single population ignores linkage disequilibrium \citep{Camerlenghi2021, Masoero2022,chakraborty2019using, zou2016quantifying,gravel2014predicting,ionita2009estimating}.
Likewise, we will ignore any potential spatial relationship between the $\genericvariantcount$ across $\variantidx$.
By the first (exchangeability) assumption, we have that $\genericvariantcount$ is an exchangeable sequence across $\unitidx$
for any fixed $\popidx$ and $\variantidx$.
We use the second assumption (ignoring linkage disequilibrium) to, momentarily, treat each sequence $(\genericvariantcount)_\unitidx$ separately.
Then de Finetti's theorem \citep{de1929funzione} applied to any $(\genericvariantcount)_\unitidx$
implies the existence of an unknown parameter $\genericvariantfreq \in [0,1]$
representing the frequency of variant $\variantidx$ in population $\popidx$,
such that $\genericvariantcount \sim \bernoullirv{\genericvariantfreq}$ i.i.d.\ across $\unitidx$.}

\revision{More precisely, we will collect 
variant frequencies across populations for the $\variantidx$-th variant 
in the vector $\genericvecvariantfreq = (\variantfreq{1}{\variantidx}, \variantfreq{2}{\variantidx})$. 
Then we can match the frequency vector $\genericvecvariantfreq$ with its variant label in a vector-valued 
measure: $\randommeas :=  \sum_{\variantidx=1}^{\infty} \genericvecvariantfreq \dirac{\genericvariantlabel}$.
We assume $\genericvariantcount \sim \bernoullirv{\genericvariantfreq}$ i.i.d.\ across $\unitidx$
and independent across $\popidx$ and $\variantidx$, conditional on $\randommeas$.
When we draw the $\genericvariantcount$ conditional on 
$\randommeas$ in this way and collect them in the measures 
$\pilotdata$, we say that $\pilotdata$
is drawn according to a \emph{multiple-population Bernoulli process} ($\mBeP$) with measure parameter $\randommeas$ and count-vector parameter $\vecsamplesize$.}

It remains to specify a prior over the frequencies.
In the one-population case, \citet{Masoero2022} generate the frequencies according
to a Poisson point process (PPP) with a diffuse (i.e., non-atomic) rate measure. We take an analogous
approach. We treat the labels $\genericvariantlabel$
as values of convenience; without loss of generality, 
we can imagine them as generated i.i.d.\ uniform on $[0,1]$ in what follows, to ensure almost sure 
uniqueness.
Then we write $\randommeas \sim \PPP(\levync)$ to indicate that $\randommeas$
is generated by drawing $\{\genericvecvariantfreq\}_{\variantidx=1}^{\infty}$ from a PPP with rate measure $\levync(\de\vecvariantfreq{})$ and drawing the $\genericvariantlabel$
i.i.d.\ uniform on $[0,1]$.

We call the resulting model the \emph{two-population model}. The two population model
can be written succinctly as follows, where it remains to choose a rate measure $\levync(\de\vecvariantfreq{})$:
\begin{equation}
		\textrm{ prior } \randommeas \sim \PPP(\levync) \quad \quad \quad
		\textrm{ and likelihood } \pilotdata \mid \randommeas \sim \mBeP(\randommeas, \vecsamplesize).
		\label{eq:generative_model}
\end{equation}

\revision{
Before continuing, we observe that one natural approach to constructing a prior over the variant frequencies
could employ a hierarchical model \citep[e.g.,][]{thibaux2007hierarchical, masoero2018posterior, james2023bayesian}.
Unlike a hierarchical model, we find that the main practical benefit of putting a Poisson point process rate measure on $\vecvariantfreq{} \in [0,1]^2$ (which is ultimately our approach in \cref{sec:our_model}) is that it leads to an efficient scheme for posterior inference, and in particular avoids the need to perform any high-dimensional integral.
See Section S2.1 in the Appendix for a discussion of concerns that arise in a hierarchical modeling approach.
However, we also find that some care must be taken in the choice of $\levync$. Before describing our model in \cref{sec:our_model},
we first rule out a simpler choice of $\levync$ in \cref{sec:cant_have_it_all}.
}



